% !PRINTED 518
\weberpage{518}
%<<CONT
einstimmung mit der Anzahl der unabhängigen Charaktere. Auf ähnliche Art ergeben sich die übrigen Fälle von (16).
% SOURCE-ORIGINAL-SIC: The registered page cites (16) for the remaining cases, although the preceding four-case adjunction table is numbered (17); the printed reference is preserved.

In den noch übrigen Fällen, nämlich $m \equiv 12\ (\mathrm{mod}\ 16)$ und $m \equiv 0\ (\mathrm{mod}\ 8)$, ist die vorstehende Zerlegung zwar auch noch anwendbar, in den so gewonnenen Teilgleichungen sind aber immer noch je zwei oder je vier Geschlechter vereinigt, entsprechend den Charakteren
\[
\left(\frac{-1}{A}\right), \ \left(\frac{2}{A}\right).
\]
Wir leiten also noch eine zweite Transformationsformel wie (8) her, indem wir die Gleichung
\[
4n = y^2 - Dx^2
\]
(8), § 136, folgendermaßen befriedigen:
\[
m = m'm'' \equiv 0\ (\mathrm{mod}\ 4),
\]
\[
x = 2, \quad y = 2(m' - m''), \quad n = (m'+m'')^2,
\]
worin wieder $m''$ ungerade und durch kein Quadrat teilbar angenommen ist, aber auch $= 1$ sein kann, und aus § 136, (9) erhält man
\[
\alpha = B + m' - m'', \quad \beta = 2A, \quad \gamma \equiv -2C\ (\mathrm{mod}\ 8),
\]
\[
\alpha m'' = AC - \left(\tfrac12 B - m''\right)^2,
\]
\[
\sqrt{\alpha + \beta\omega} = \sqrt{m'} + i\sqrt{m''}.
\]
Nehmen wir der Einfachheit halber $B \equiv 0\ (\mathrm{mod}\ 8)$, was nötigenfalls durch Übergang zu einer parallelen Form erreicht wird, so ergibt sich $m' \equiv C\ (\mathrm{mod}\ 8)$, $\alpha \equiv C - m''\ (\mathrm{mod}\ 8)$ und folglich
\[
i^{3\frac{\alpha-1}{2}} e^{\frac{\pi i}{4}\alpha(\gamma-\beta)} = (-1)^{\frac14 C} (-1)^{\frac{A+1}{2}} i^{3\frac{m''-1}{2}},
\]
\[
\left(\frac{\beta}{\alpha}\right) = \left(\frac{2}{m''}\right)\left(\frac{A}{m''}\right)\left(\frac{2}{\alpha m''}\right)\left(\frac{A}{\alpha m''}\right).
\]
\[
\left(\frac{A}{\alpha m''}\right) = 1, \quad \left(\frac{2}{\alpha m''}\right) = (-1)^{\frac14 C},
\]
und daraus ergibt sich nach § 136, (12), (13):
\weq{18}{M = -\left(\frac{2}{m''}\right)\left(\frac{A}{m''}\right)\left(\frac{-1}{A}\right) \left(i^{\frac{m''-1}{2}}\sqrt{m'} + i^{\frac{m''+1}{2}}\sqrt{m''}\right)^3.}
Diese Formel ergänzt die vorige für den Fall $m \equiv 12\ (\mathrm{mod}\ 16)$, und zeigt, daß auch in diesem Falle die vollständige Zerfällung der Klassengleichung durch Adjunktion von
\[
\sqrt{p}, \ \sqrt{p'}, \ \sqrt{p''}, \ \ldots, \ \sqrt{q}, \ \sqrt{q'}, \ \sqrt{q''}, \ \ldots
\]
geschieht.
